\documentclass[a4paper,oneside,12pt]{report}
% Package yang diperlukan
\usepackage[table]{xcolor} % Untuk menggunakan warna
\usepackage{graphicx}      % Untuk memasukkan gambar
\usepackage{titling}       % Untuk kustomisasi halaman judul
\usepackage{geometry}      % Untuk mengatur margin halaman
\usepackage{listings}      % Untuk menampilkan blok kode
\usepackage{float}         % Untuk kontrol penempatan float
\usepackage{hyperref}      % Untuk membuat hyperlink
\usepackage{indentfirst}   % Untuk memberi indentasi pada paragraf pertama
\usepackage{url}           % Untuk format URL
\usepackage{fancyhdr}      % Untuk kustomisasi header dan footer
\usepackage{tocloft}       % Untuk kustomisasi Daftar Isi
\usepackage{bookmark}      % Untuk manajemen bookmark PDF
\usepackage{amsmath}       % Untuk notasi matematika
\usepackage[most]{tcolorbox} % Untuk box output terminal

% Paket untuk tabel
\usepackage{array}
\usepackage{booktabs}
\usepackage{tabularx}
\usepackage{multirow}
\usepackage{colortbl}

% Macro dan warna untuk tabel
\newcommand{\thead}[1]{\textcolor{white}{\textbf{#1}}}
\definecolor{tableheader}{RGB}{30, 58, 138}
\definecolor{tablezebra}{RGB}{248, 250, 252}
\definecolor{tableborder}{RGB}{203, 213, 225}

%======================================================================
% PENGATURAN GEOMETRI HALAMAN
%======================================================================
\geometry{left=3cm, right=2.5cm, top=2cm, bottom=2.5cm}
\sloppy

%======================================================================
% PENGATURAN WARNA UNTUK BLOK KODE
%======================================================================
\definecolor{codegreen}{rgb}{0,0.6,0}
\definecolor{codegray}{rgb}{0.5,0.5,0.5}
\definecolor{codepurple}{rgb}{0.58,0,0.82}
\definecolor{backcolour}{rgb}{0.95,0.95,0.92}
\definecolor{darkblue}{rgb}{0,0,0.6}
\definecolor{outputbg}{RGB}{241,245,249}
\definecolor{outputborder}{RGB}{203,213,225}

%======================================================================
% PENGATURAN BLOK KODE (LISTINGS)
%======================================================================
\lstdefinestyle{commonstyle}{
    backgroundcolor=\color{backcolour},
    breakatwhitespace=false,
    breaklines=true,
    captionpos=b,
    keepspaces=true,
    numbers=left,
    numbersep=5pt,
    showspaces=false,
    showstringspaces=false,
    showtabs=false,
    tabsize=2,
    frame=single,
    basicstyle=\ttfamily\footnotesize,
    numberstyle=\tiny\color{codegray},
    columns=flexible,
    postbreak=\mbox{\textcolor{red}{$\hookrightarrow$}\space},
    breakindent=0pt,
    breakautoindent=false
}

\lstdefinestyle{pythonstyle}{
    style=commonstyle,
    language=Python,
    commentstyle=\color{codegreen},
    keywordstyle=\color{darkblue},
    stringstyle=\color{codepurple},
    emph={import,from,class,def,for,while,if,else,elif,try,except,finally,as,lambda,return},
    emphstyle=\color{darkblue}
}

% Macro untuk screenshot gambar Colab
\newcommand{\screenshotimage}[4][0.88\textwidth]{%
    \begin{figure}[H]
        \centering
        \includegraphics[width=#1]{#2}
        \caption{#3}
        \label{fig:#4}
    \end{figure}
}

%======================================================================
% PENGATURAN LAIN-LAIN
%======================================================================
\renewcommand{\figurename}{Gambar}
\renewcommand{\thefigure}{\arabic{figure}}
\renewcommand{\tablename}{Tabel}

\hypersetup{
    colorlinks=true,
    linkcolor=black,
    filecolor=black,
    urlcolor=black,
    citecolor=black,
    pdftitle={Laporan Praktikum Penambangan Data Pertemuan VI},
    pdfpagemode=FullScreen,
    breaklinks=true
}

\pagestyle{fancy}
\fancyhf{}
\fancyfoot[C]{\thepage}
\renewcommand{\headrulewidth}{0pt}

\fancypagestyle{plain}{
    \fancyhf{}
    \fancyfoot[C]{\thepage}
    \renewcommand{\headrulewidth}{0pt}
}

%======================================================================
% PENGATURAN STRUKTUR BAB DAN DAFTAR ISI
%======================================================================
\renewcommand{\contentsname}{Daftar Isi}
\renewcommand{\bibname}{Daftar Pustaka}

\renewcommand{\thechapter}{}
\renewcommand{\chaptername}{}

\renewcommand{\thesection}{\arabic{chapter}.\arabic{section}}
\renewcommand{\thesubsection}{\thesection.\arabic{subsection}}

\setcounter{secnumdepth}{2}
\setcounter{tocdepth}{2}

\renewcommand{\cftchapdotsep}{\cftdotsep}
\renewcommand{\cftchapleader}{\cftdotfill{\cftchapdotsep}}
\setlength{\cftchapindent}{0em}
\setlength{\cftsecindent}{2em}
\setlength{\cftsubsecindent}{4em}
\setlength{\cftchapnumwidth}{0em}
\setlength{\cftsecnumwidth}{2.5em}
\setlength{\cftsubsecnumwidth}{3.2em}
\renewcommand{\cftchappresnum}{}
\renewcommand{\cftchapaftersnum}{}

%======================================================================
% DOKUMEN UTAMA
%======================================================================
\begin{document}

%======================================================================
% HALAMAN JUDUL (COVER)
%======================================================================
\newgeometry{left=3cm, right=2.5cm, top=2.5cm, bottom=3cm}
\begin{titlingpage}
\begin{center}
\vspace*{0.5cm}
{\large \textbf{\MakeUppercase{Laporan Praktikum}}} \\[0.4cm]
{\large \textbf{\MakeUppercase{Praktikum Penambangan Data}}} \\[0.4cm]
{\large \textbf{\MakeUppercase{Pertemuan VI}}} \\[0.4cm]
{\large \textbf{\MakeUppercase{Clustering}}} \\[1.5cm]

\includegraphics[height=7cm]{lambang ugm.png} \\[1cm]

{\normalsize \textbf{Disusun oleh:}} \\[0.4cm]
\begin{tabular}{@{}l l @{}}
  \textbf{Nama}           & : Januarsyah Akbar \\
  \textbf{NIM}            & : 24/535846/SV/24314 \\
  \textbf{Kelas}          & : B2 \\
  \textbf{Dosen Pengampu} & : Dr. Ganjar Alfian, S.T., M.Eng. \\
                          & \phantom{:} Dr. Imam Fahrurrozi, S.T., M.Cs.
\end{tabular} \\[1.2cm]

{\normalsize \textbf{PROGRAM STUDI D-IV}} \\[0.2cm]
{\normalsize \textbf{TEKNOLOGI REKAYASA PERANGKAT LUNAK}} \\[0.2cm]
{\normalsize \textbf{DEPARTEMEN TEKNIK ELEKTRO DAN INFORMATIKA}} \\[0.2cm]
{\normalsize \textbf{SEKOLAH VOKASI}} \\[0.2cm]
{\normalsize \textbf{UNIVERSITAS GADJAH MADA}} \\[0.2cm]
{\normalsize \textbf{YOGYAKARTA}} \\[0.2cm]
{\normalsize \textbf{2026}} \\[1.2cm]
\end{center}
\end{titlingpage}
\restoregeometry

\thispagestyle{empty}
\clearpage
\stepcounter{page}
\thispagestyle{empty}

% --- Daftar Isi ---
\addcontentsline{toc}{chapter}{Daftar Isi}
\tableofcontents
\clearpage

\pagenumbering{arabic}
\setcounter{page}{3}

%======================================================================
% BAB I: TUJUAN PRAKTIKUM
%======================================================================
\chapter[Tujuan Praktikum]{Tujuan Praktikum}

\begin{enumerate}
    \item Mahasiswa memahami konsep dasar \textit{unsupervised learning} dan perannya dalam pembelajaran mesin.
    \item Mahasiswa mampu menerapkan berbagai teknik \textit{clustering} seperti K-Means, K-Medoids, dan DBSCAN untuk mengelompokkan data berdasarkan karakteristik tertentu.
    \item Mahasiswa dapat membandingkan kualitas hasil \textit{clustering} menggunakan metrik evaluasi seperti Inertia dan Silhouette Score.
    \item Mahasiswa memahami konsep \textit{Association Rule} dan mampu menerapkannya untuk menemukan pola hubungan antar-item dalam dataset.
\end{enumerate}
\clearpage

%======================================================================
% BAB II: DASAR TEORI
%======================================================================
\chapter[Dasar Teori]{Dasar Teori}

\section{Unsupervised Learning}
\textit{Unsupervised Learning} merupakan salah satu paradigma pembelajaran mesin di mana model dilatih untuk menemukan struktur tersembunyi atau pola dari sekumpulan data tanpa bantuan label target eksternal. Algoritma ini menganalisis hubungan antar-variabel data tanpa intervensi atau arahan manusia. Fokus utamanya adalah eksplorasi data deskriptif dan pengelompokan menggunakan algoritma seperti K-Means, Artificial Neural Network (ANN), dan Gaussian Mixture Model (GMM) yang menawarkan fleksibilitas tinggi dalam menjelajahi karakteristik data yang belum diketahui sebelumnya (Nurhalizah \& Ardianto, 2024).

Keunggulan utama \textit{unsupervised learning} terletak pada kemampuannya mengolah data mentah dalam jumlah besar yang belum ber-label. Teknik ini sangat populer pada aplikasi deteksi anomali, rekomendasi produk, serta analisis segmentasi pasar di berbagai bidang industri (Wijoyo et al., 2024).

\section{Clustering}
\textit{Clustering} atau klasterisasi adalah teknik analisis data statistik yang bertujuan mengelompokkan sekumpulan objek atau sampel data ke dalam beberapa kelompok (\textit{cluster}). Setiap kelompok berisi objek yang memiliki tingkat kemiripan (\textit{homogeneity}) tinggi satu sama lain, sementara antar-kelompok memiliki tingkat perbedaan (\textit{heterogeneity}) yang maksimal. Analisis klaster telah dimanfaatkan secara luas di bidang statistik, ekonomi, kesehatan, hingga sosiologi kependudukan (Afidaha \& Masrukan, 2023).

\section{Algoritma K-Means}
K-Means merupakan salah satu algoritma \textit{partitioning clustering} non-hirarki yang paling populer. Algoritma ini mempartisi $N$ sampel data ke dalam $K$ buah klaster berdasarkan jarak Euclidean terdekat terhadap titik pusat massa (\textit{centroid}). Proses K-Means beriterasi secara independen dimulai dari penentuan posisi awal $K$ centroid secara acak, pengalokasian data ke centroid terdekat, hingga pembaruan lokasi centroid berdasarkan rata-rata (\textit{mean}) koordinat seluruh anggota klaster (Sulistiyawati \& Supriyanto, 2021).

Tujuan utama dari optimasi K-Means adalah meminimalkan fungsi objektif \textit{Within-Cluster Sum of Squares} (WCSS) atau Inertia. Meskipun cepat dan efisien untuk dataset berskala besar, K-Means rentan terhadap penentuan centroid awal dan sangat sensitif terhadap keberadaan pencilan (\textit{outliers}).

\section{Algoritma K-Medoids}
K-Medoids atau \textit{Partitioning Around Medoids} (PAM) merupakan varian algoritma pemartisian yang menggunakan titik data aktual (\textit{medoid}) sebagai pusat klaster, bukan nilai rata-rata imajiner seperti pada K-Means. Pendekatan ini membuat K-Medoids jauh lebih resisten dan tangguh (\textit{robust}) terhadap pengaruh \textit{noise} maupun pencilan ekstrem pada data (Herviany et al., 2021).

Meskipun secara komputasi K-Medoids memerlukan iterasi perhitungan matriks jarak yang lebih kompleks dibanding K-Means, algoritma ini sangat cocok digunakan pada data yang memiliki pencilan atau ketika pusat klaster dituntut merupakan sampel data yang riil.

\section{Algoritma DBSCAN}
\textit{Density-Based Spatial Clustering of Applications with Noise} (DBSCAN) adalah algoritma \textit{clustering} berbasis kepadatan area. Berbeda dari K-Means dan K-Medoids yang mengharuskan penentuan jumlah klaster $K$ di awal, DBSCAN menentukan klaster berdasarkan kerapatan titik data yang dibatasi oleh radius epsilon ($\epsilon$) dan jumlah minimum tetangga (\textit{minPts}). Algoritma ini sangat unggul dalam menemukan klaster berbentuk arbitrer serta mampu mengisolasi titik pencilan sebagai \textit{noise} secara otomatis (Ester et al., 1996; Han et al., 2011).

\section{Association Rule Mining}
\textit{Association Rule Mining} merupakan teknik penambangan data untuk menemukan aturan keterkaitan atau pola asosiasi antar-item dalam sekumpulan transaksi. Aturan umumnya dinyatakan dalam bentuk $A \rightarrow B$, yang berarti kemunculan item $A$ cenderung diikuti oleh kemunculan item $B$. Evaluasi kualitas aturan asosiasi diukur menggunakan tiga metrik utama, yaitu \textit{support} (frekuensi kemunculan), \textit{confidence} (kepastian probabilitas bersyarat), dan \textit{lift} (kekuatan keterikatan aturan dibanding acak) (Agrawal et al., 1993; Han et al., 2011).

\section{Inertia (Within-Cluster Sum of Squares)}
Inertia atau WCSS adalah metrik evaluasi internal pada K-Means yang mengukur tingkat kerapatan (\textit{compactness}) di dalam klaster. Nilai Inertia dihitung sebagai jumlah total kuadrat jarak antara setiap titik sampel terhadap centroid klaster masing-masing. Semakin kecil nilai Inertia, semakin rapat dan representatif anggota di dalam klaster tersebut. Metrik Inertia menjadi acuan utama pada metode siku (\textit{Elbow Method}) untuk menentukan jumlah klaster optimal (Bishop, 2006).

\section{Silhouette Score}
\textit{Silhouette Score} merupakan metrik evaluasi kualitas pengelompokan yang mengukur seberapa baik suatu titik data ditempatkan pada klasternya dibandingkan dengan klaster tetangga terdekatnya. Nilai Silhouette Score berada pada rentang -1 hingga +1. Nilai yang mendekati +1 mengindikasikan pemisahan antar-klaster yang sangat jelas dan baik, nilai 0 menunjukkan adanya tumpang tindih (\textit{overlap}) antar-klaster, sedangkan nilai negatif menandakan adanya kesalahan pengelompokan (Rousseeuw, 1987; Tan et al., 2019).

\clearpage

%======================================================================
% BAB III: HASIL DAN PEMBAHASAN
%======================================================================
\chapter[Hasil dan Pembahasan]{Hasil dan Pembahasan}

\section{Langkah Praktikum}

\subsection{Langkah Percobaan Clustering Credit Card Dataset}

Praktikum dimulai dengan mempersiapkan lingkungan pengembangan di Google Colab. Seluruh pustaka Python yang diperlukan untuk manipulasi data, kalkulasi statistik, standarisasi, pemodelan kluster, dan evaluasi diimpor ke dalam notebook.

\screenshotimage{gambar/colab-p6-imports.png}{Import Pustaka dan Setup Dependensi Google Colab}{colab-imports}

Setelah dependensi siap, dataset perilaku nasabah kartu kredit (\textit{Credit Card Dataset for Clustering}) dimuat ke dalam DataFrame pandas. Dataset memiliki dimensi sebanyak 8.950 baris sampel nasabah dan 18 atribut transaksi.

\screenshotimage{gambar/colab-p6-read-dataset.png}{Cuplikan Lima Baris Pertama Credit Card Dataset (df.head())}{colab-read}

Pengecekan ukuran matriks dataset mengonfirmasi keberadaan 8.950 sampel nasabah dengan 18 kolom fitur secara utuh.

\screenshotimage[0.55\textwidth]{gambar/colab-p6-shape.png}{Dimensi dan Ukuran Matriks Credit Card Dataset (df.shape)}{colab-shape}

Pemeriksaan ringkasan statistik deskriptif dilakukan untuk memahami distribusi min, max, mean, dan standar deviasi pada setiap atribut fitur transaksi nasabah.

\screenshotimage{gambar/colab-p6-describe.png}{Ringkasan Statistik Deskriptif Atribut Dataset (df.describe().T)}{colab-describe}

Selanjutnya dilakukan analisis korelasi Pearson untuk mengidentifikasi tingkat keterkaitan linier antar-fitur transaksi, baik melalui tabel matriks korelasi numerik maupun visualisasi peta panas (\textit{heatmap}).

\screenshotimage{gambar/colab-p6-corr-table.png}{Tabel Matriks Korelasi Linier Antar-Atribut Numerik (df.corr())}{colab-corr-table}

\screenshotimage{gambar/colab-p6-corr.png}{Heatmap Korelasi Pearson Antar-Fitur Transaksi Credit Card}{colab-corr}

Tahap pembersihan data (\textit{data cleaning}) diawali dengan menghapus atribut \texttt{CUST\_ID} karena atribut tersebut hanya berfungsi sebagai identitas nasabah dan tidak memiliki relevansi terhadap proses pengelompokan.

\screenshotimage{gambar/colab-p6-drop-custid.png}{Penghapusan Atribut Identitas CUST\_ID dari DataFrame}{colab-drop-custid}

Pemeriksaan nilai yang hilang (\textit{missing values}) menunjukkan adanya nilai kosong pada atribut \texttt{MINIMUM\_PAYMENTS} (313 entri) dan \texttt{CREDIT\_LIMIT} (1 entri). Nilai kosong tersebut diimputasi menggunakan nilai median dari masing-masing atribut.

\screenshotimage{gambar/colab-p6-isnull.png}{Pengecekan Missing Values pada Dataset Sebelum Imputasi}{colab-isnull}

\screenshotimage{gambar/colab-p6-imputation.png}{Proses Imputasi Missing Values Menggunakan Nilai Median}{colab-imputation}

Sebelum menjalankan algoritma kluster, seluruh fitur distandarisasi menggunakan \texttt{StandardScaler}. Standarisasi ini sangat krusial agar fitur berskala besar seperti \texttt{BALANCE} atau \texttt{PURCHASES} tidak mendominasi perhitungan jarak Euclidean terhadap fitur rasio atau frekuensi.

\screenshotimage{gambar/colab-p6-scaling.png}{Matriks Hasil Standarisasi Fitur Menggunakan StandardScaler}{colab-scaling}

Pengujian jumlah kluster optimal $k=1 \dots 10$ dilakukan menggunakan metode siku (\textit{Elbow Method}) pada algoritma K-Means. Penurunan nilai Inertia dicatat pada setiap iterasi.

\screenshotimage{gambar/colab-p6-kmeans-elbow-data.png}{Daftar Nilai Inertia K-Means untuk k=1 hingga k=10}{colab-km-elbow-data}

\screenshotimage{gambar/colab-p6-kmeans-elbow-plot.png}{Grafik Kurva Elbow Method pada Algoritma K-Means}{colab-km-elbow-plot}

Untuk menentukan titik tekukan siku (\textit{knee point}) secara objektif dan akurat, digunakan pustaka \texttt{KneeLocator} yang menghasilkan titik siku optimal pada $k=4$.

\screenshotimage{gambar/colab-p6-kmeans-knee-plot.png}{Deteksi Siku Optimal K-Means (k=4) Menggunakan KneeLocator}{colab-km-knee-plot}

Proses serupa dijalankan untuk menguji nilai Inertia pada algoritma K-Medoids untuk rentang kluster $k=1 \dots 10$.

\screenshotimage{gambar/colab-p6-kmedoids-elbow-data.png}{Daftar Nilai Inertia K-Medoids untuk k=1 hingga k=10}{colab-kmed-elbow-data}

\screenshotimage{gambar/colab-p6-kmedoids-elbow-plot.png}{Grafik Kurva Elbow Method pada Algoritma K-Medoids}{colab-kmed-elbow-plot}

\screenshotimage{gambar/colab-p6-kmedoids-knee-plot.png}{Deteksi Siku Optimal K-Medoids (k=4) Menggunakan KneeLocator}{colab-kmed-knee-plot}

Selain metode siku, evaluasi kualitas pemisahan kluster diuji menggunakan koefisien siluet (\textit{Silhouette Score}). Pada algoritma K-Means, skor puncak tertinggi diperoleh pada jumlah kluster $k=2$ dengan perolehan skor sebesar 0,2795.

\screenshotimage{gambar/colab-p6-kmeans-sh-data.png}{Output Perhitungan Silhouette Score K-Means}{colab-km-sh-data}

\screenshotimage{gambar/colab-p6-kmeans-sh-plot.png}{Grafik Silhouette Score K-Means Memuncak Tajam pada k=2}{colab-km-sh-plot}

Pengujian koefisien siluet pada K-Medoids juga menunjukkan nilai tertinggi berada pada $k=2$ dengan skor 0,1945, yang kemudian menurun seiring bertambahnya jumlah kluster.

\screenshotimage{gambar/colab-p6-kmedoids-sh-data.png}{Output Perhitungan Silhouette Score K-Medoids}{colab-kmed-sh-data}

\screenshotimage{gambar/colab-p6-kmedoids-sh-plot.png}{Grafik Silhouette Score K-Medoids terhadap Jumlah Kluster}{colab-kmed-sh-plot}

Model K-Means kemudian dilatih menggunakan parameter optimal $k=2$. Label hasil pengelompokan ditambahkan ke dalam DataFrame, dan titik pusat massa (\textit{centroid}) kluster diinspeksi koordinatnya.

\screenshotimage{gambar/colab-p6-kmeans-train.png}{Pelatihan Model K-Means (k=2) dan Penambahan Label Kluster}{colab-km-train}

\screenshotimage{gambar/colab-p6-kmeans-centroids.png}{Dimensi dan Nilai Koordinat Centroid K-Means (2 Centroid x 17 Fitur)}{colab-km-centroids}

Persebaran data divisualisasikan dalam bentuk diagram sebar (\textit{scatter plot}) 2D antara fitur \texttt{PURCHASES} dan \texttt{PAYMENTS}.

\screenshotimage{gambar/colab-p6-kmeans-scatter2d.png}{Visualisasi 2D Scatter Plot K-Means Clustering (k=2)}{colab-km-scatter2d}

Untuk memperoleh gambaran sebaran data yang lebih kaya, dibuat pula visualisasi 3D yang melibatkan tiga atribut utama, yaitu \texttt{PURCHASES}, \texttt{PAYMENTS}, dan \texttt{BALANCE}.

\screenshotimage{gambar/colab-p6-kmeans-scatter3d.png}{Visualisasi 3D Scatter Plot K-Means (PURCHASES vs PAYMENTS vs BALANCE)}{colab-km-scatter3d}

Model K-Medoids dilatih menggunakan $k=4$ berdasarkan hasil metode siku untuk melihat segmentasi nasabah secara lebih mendalam ke dalam empat kelompok perilaku. Titik sampel data riil yang menjadi medoid kluster diekstraksi koordinatnya.

\screenshotimage{gambar/colab-p6-kmedoids-train.png}{Pelatihan Model K-Medoids (k=4) dan Penambahan Label Kluster}{colab-kmed-train}

\screenshotimage{gambar/colab-p6-kmedoids-medoids.png}{Dimensi dan Nilai Koordinat Medoids Aktual K-Medoids (4 Medoid x 17 Fitur)}{colab-kmed-medoids}

\screenshotimage{gambar/colab-p6-kmedoids-scatter2d.png}{Visualisasi 2D Scatter Plot K-Medoids Clustering (k=4)}{colab-kmed-scatter2d}

Analisis profil keempat kluster K-Medoids dilakukan dengan membandingkan nilai rata-rata atribut \texttt{PURCHASES}, \texttt{PAYMENTS}, dan \texttt{BALANCE} menggunakan diagram batang.

\screenshotimage{gambar/colab-p6-cluster-analysis-barplot.png}{Bar Plot Perbandingan Rata-rata Fitur Utama per Kluster K-Medoids}{colab-bar-analysis}

Berdasarkan hasil visualisasi dan statistik kluster K-Medoids, profil segmen nasabah dirumuskan sebagai berikut:
\begin{enumerate}
    \item \textbf{Cluster 0 (Nasabah Hemat atau Pasif)}: Memiliki rata-rata \texttt{PURCHASES}, \texttt{PAYMENTS}, dan \texttt{BALANCE} paling rendah di antara seluruh kelompok nasabah.
    \item \textbf{Cluster 1 (Nasabah Reguler)}: Memiliki nilai transaksi dan saldo pada tingkat moderat.
    \item \textbf{Cluster 2 (Nasabah Transaktif Tinggi)}: Memiliki rata-rata \texttt{PURCHASES} dan \texttt{PAYMENTS} tertinggi, mencirikan nasabah yang sangat aktif berbelanja dan tertib melunasi tagihan.
    \item \textbf{Cluster 3 (Nasabah Peminjam Kas atau Pembawa Saldo)}: Memiliki nilai saldo \texttt{BALANCE} paling tinggi dengan frekuensi transaksi belanja langsung yang sedang.
\end{enumerate}

Sebagai penutup langkah percobaan, ditunjukkan demonstrasi penggunaan \texttt{LabelEncoder} dari modul \texttt{sklearn.preprocessing} untuk mengonversi variabel bertipe kategorial menjadi representasi numerik.

\screenshotimage{gambar/colab-p6-label-encoder.png}{Demonstrasi Penerapan LabelEncoder pada Kolom Kategorial}{colab-label-encoder}

\section{Tugas}

\subsection*{Tugas 1: Preparation \& Analisis Dataset Credit Card (CCData)}
\addcontentsline{toc}{subsection}{Tugas 1: Preparation \& Analisis Dataset Credit Card (CCData)}

Dataset \texttt{CC GENERAL.csv} atau dikenal sebagai Kaggle Credit Card Dataset (CCData) merupakan kumpulan data publik yang merekam aktivitas transaksi 8.950 nasabah kartu kredit aktif selama periode observasi 6 bulan.

Tujuan utama analisis klaster pada dataset ini mencakup:
\begin{enumerate}
    \item \textbf{Segmentasi Nasabah}: Membagi basis pelanggan ke dalam beberapa kluster homogen berdasarkan kebiasaan belanja, penarikan tunai di ATM, dan kedisiplinan pelunasan tagihan.
    \item \textbf{Pengembangan Strategi Pemasaran Tertarget}: Membantu institusi perbankan merancang program promosi yang tepat sasaran, seperti pemberian poin penghargaan bagi nasabah transaktif atau fasilitas cicilan fleksibel bagi nasabah pembawa saldo.
    \item \textbf{Manajemen Risiko Kredit}: Mengidentifikasi kelompok pengguna kartu kredit yang rentan mengalami gagal bayar melalui pola akumulasi saldo tinggi yang dibarengi tingkat pembayaran tagihan minimum yang rendah.
\end{enumerate}

Rincian 18 atribut fitur yang tercakup di dalam dataset CCData disajikan pada Tabel \ref{tab:cc-features}.

\begin{table}[H]
\centering
\caption{Deskripsi Rincian 18 Fitur Dataset Credit Card (CCData)}
\label{tab:cc-features}
\renewcommand{\arraystretch}{1.25}
\rowcolors{2}{tablezebra}{white}
\small
\begin{tabularx}{\textwidth}{>{\bfseries}l l X}
\toprule
\rowcolor{tableheader}
\thead{Nama Fitur} & \thead{Tipe Data} & \thead{Deskripsi dan Makna Bisnis} \\
\midrule
\texttt{CUST\_ID} & Kategorial & Identitas unik bagi setiap pemilik kartu kredit. \\
\texttt{BALANCE} & Kontinu & Sisa saldo tagihan yang belum dilunasi oleh nasabah. \\
\texttt{BALANCE\_FREQUENCY} & Kontinu & Rasio frekuensi pembaruan saldo (skala 0 hingga 1). \\
\texttt{PURCHASES} & Kontinu & Total akumulasi nominal transaksi pembelian. \\
\texttt{ONEOFF\_PURCHASES} & Kontinu & Nominal transaksi pembelian sekali bayar (tanpa cicilan). \\
\texttt{INSTALLMENTS\_PURCHASES} & Kontinu & Total nominal transaksi pembelian berbasis cicilan. \\
\texttt{CASH\_ADVANCE} & Kontinu & Total nominal penarikan uang tunai di muka dari ATM. \\
\texttt{PURCHASES\_FREQUENCY} & Kontinu & Frekuensi keaktifan transaksi pembelian (skala 0 hingga 1). \\
\texttt{ONEOFF\_PURCHASES\_FREQ} & Kontinu & Frekuensi transaksi pembelian sekali bayar. \\
\texttt{PURCHASES\_INSTALL\_FREQ} & Kontinu & Frekuensi transaksi pembelian berbasis cicilan. \\
\texttt{CASH\_ADVANCE\_FREQ} & Kontinu & Frekuensi penarikan tunai uang di muka. \\
\texttt{CASH\_ADVANCE\_TRX} & Diskret & Jumlah total frekuensi transaksi penarikan tunai. \\
\texttt{PURCHASES\_TRX} & Diskret & Jumlah total frekuensi transaksi pembelian yang berhasil. \\
\texttt{CREDIT\_LIMIT} & Kontinu & Batas maksimum limit kartu kredit yang diberikan bank. \\
\texttt{PAYMENTS} & Kontinu & Total jumlah pembayaran tagihan yang telah disetor. \\
\texttt{MINIMUM\_PAYMENTS} & Kontinu & Batas pembayaran tagihan minimum yang diwajibkan. \\
\texttt{PRC\_FULL\_PAYMENT} & Kontinu & Persentase pelunasan tagihan secara penuh oleh nasabah. \\
\texttt{TENURE} & Diskret & Durasi keanggotaan kartu kredit nasabah (dalam bulan). \\
\bottomrule
\end{tabularx}
\end{table}

\subsection*{Tugas 2: Clustering Dataset Air Traffic Passengers Statistics}
\addcontentsline{toc}{subsection}{Tugas 2: Clustering Dataset Air Traffic Passengers Statistics}

Tugas kedua mengaplikasikan algoritma K-Means dan K-Medoids pada dataset \textit{Air Traffic Passengers Statistics} yang disediakan oleh Pemerintah Kota San Francisco (SF Open Data). Dataset ini merekam data statistik lalu lintas penumpang penerbangan di Bandara Internasional San Francisco (SFO).

\screenshotimage{gambar/colab-p6-air-head.png}{Cuplikan Data Statistik Penumpang Penerbangan SFO (df\_air.head())}{colab-air-head}

Struktur data diperiksa untuk memastikan jumlah sampel baris dan keberadaan nilai kosong. Dataset terdiri dari 1.000 sampel data transaksi penerbangan dengan atribut utama seperti \texttt{operating\_airline}, \texttt{geo\_summary}, \texttt{geo\_region}, \texttt{activity\_type\_code}, dan \texttt{passenger\_count}.

\screenshotimage{gambar/colab-p6-air-info.png}{Ringkasan Struktur Informasi dan Pengecekan Missing Values Dataset Air Traffic}{colab-air-info}

Tahap pra-pemrosesan data dilakukan dengan menyaring fitur numerik dan mengodekan atribut kategorial kuncinya (\texttt{operating\_airline}, \texttt{geo\_summary}, \texttt{geo\_region}, \texttt{activity\_type\_code}, \texttt{price\_category\_code}) menjadi bentuk angka numerik menggunakan \texttt{LabelEncoder}. Selanjutnya dilakukan standarisasi fitur menggunakan \texttt{StandardScaler}.

\screenshotimage{gambar/colab-p6-air-prep.png}{Pra-pemrosesan Data, Encoding Variabel Kategorial, dan Standarisasi Fitur Air Traffic}{colab-air-prep}

Evaluasi penentuan jumlah kluster optimal dijalankan pada rentang $k=1 \dots 10$ menggunakan metode siku (\textit{Elbow Method}) dan koefisien siluet (\textit{Silhouette Score}) untuk kedua algoritma. Pustaka \texttt{KneeLocator} mendeteksi titik tekukan siku kurva Inertia pada masing-masing model.

\screenshotimage[0.6\textwidth]{gambar/colab-p6-air-optimal-clusters.png}{Penentuan Titik Siku Optimal Air Traffic Menggunakan KneeLocator}{colab-air-optimal-clusters}

\screenshotimage{gambar/colab-p6-air-elbow-sh.png}{Perbandingan Kurva Elbow dan Silhouette Score Air Traffic (K-Means vs K-Medoids)}{colab-air-elbow-sh}

Berdasarkan titik tekukan siku kurva Inertia dan skor Silhouette terbaik, ditentukan jumlah kluster optimal $k=3$. Kedua algoritma dilatih pada matriks data terstandarisasi untuk mengelompokkan karakteristik lalu lintas penumpang penerbangan.

\screenshotimage[0.7\textwidth]{gambar/colab-p6-air-eval.png}{Perbandingan Hasil Perhitungan Silhouette Score Air Traffic Clustering}{colab-air-eval}

Hasil pengujian menunjukkan bahwa model \textbf{K-Means} mencapai Silhouette Score sebesar \textbf{0,4428}, mengungguli \textbf{K-Medoids} yang memperoleh skor \textbf{0,3812}.

\screenshotimage{gambar/colab-p6-air-scatter-comp.png}{Visualisasi Scatter Plot Kluster Air Traffic (K-Means vs K-Medoids)}{colab-air-scatter-comp}

Visualisasi diagram sebar (\textit{scatter plot}) memperlihatkan bahwa K-Means mampu memisahkan titik-titik data penumpang penerbangan secara lebih rapi dan konsisten sesuai dengan tingkat kerapatan volume penumpang pada sumbu vertikal.

\screenshotimage{gambar/colab-p6-air-bar-summary.png}{Rata-rata Jumlah Penumpang per Kluster K-Means pada Dataset Air Traffic}{colab-air-bar-summary}

Berdasarkan hasil analisis distribusi nilai pada kluster K-Means, dapat dirumuskan tiga kelompok utama lalu lintas penumpang di Bandara Internasional San Francisco:
\begin{enumerate}
    \item \textbf{Kluster 0 (High-Volume Mainline Hub)}: Kelompok penerbangan dengan volume penumpang sangat tinggi ($>100.000$ penumpang per periode) yang didominasi oleh maskapai penerbangan internasional utama dan rute antar-benua.
    \item \textbf{Kluster 1 (Medium-Volume Regional Direct Flight)}: Kelompok penerbangan domestik dan regional dengan volume penumpang sedang ($20.000 - 80.000$ penumpang).
    \item \textbf{Kluster 2 (Low-Volume Niche \& Charter Flight)}: Kelompok penerbangan bervolume penumpang rendah ($<20.000$ penumpang) yang mencakup penerbangan khusus, rute kargo, carter, dan maskapai regional kecil.
\end{enumerate}

Hasil segmentasi ini memberikan implikasi manajerial yang sangat berharga bagi pihak pengelola bandara SFO. Penerbangan pada Kluster 0 memerlukan alokasi gerbang keberangkatan (\textit{boarding gate}) utama dan fasilitas klaim bagasi kapasitas besar untuk mencegah penumpukan penumpang, sedangkan Kluster 2 dapat dialokasikan pada terminal sekunder demi efisiensi operasional bandara.

Ringkasan perbandingan capaian evaluasi kuantitatif dari kedua algoritma klasterisasi pada kedua eksperimen praktikum dirangkum pada Tabel \ref{tab:clustering-comp}.

\begin{table}[H]
\centering
\caption{Ringkasan Komparasi Performa Algoritma K-Means dan K-Medoids}
\label{tab:clustering-comp}
\renewcommand{\arraystretch}{1.25}
\rowcolors{2}{tablezebra}{white}
\small
\begin{tabularx}{\textwidth}{l l c c X}
\toprule
\rowcolor{tableheader}
\thead{Dataset} & \thead{Algoritma} & \thead{$k$ Optimal} & \thead{Silhouette} & \thead{Karakteristik Hasil Kluster} \\
\midrule
\multirow{2}{*}{\textbf{Credit Card}} & K-Means & 2 & \textbf{0,2795} & Berbasis rata-rata hitung, memetakan gradasi volume transaksi. \\
 & K-Medoids & 2 & 0,1945 & Berbasis sampel riil, lebih tahan terhadap pencilan data. \\
\midrule
\multirow{2}{*}{\textbf{Air Traffic}} & K-Means & 3 & \textbf{0,4428} & Struktur bola terpartisi tegas mengikuti kepadatan penumpang. \\
 & K-Medoids & 3 & 0,3812 & Struktur stabil namun pemisahan batas kelompok kurang rapat. \\
\bottomrule
\end{tabularx}
\end{table}

\clearpage

%======================================================================
% BAB IV: KESIMPULAN
%======================================================================
\chapter[Kesimpulan]{Kesimpulan}

Praktikum pertemuan keenam memberikan pemahaman mendalam mengenai paradigma pembelajaran mesin tanpa arahan label (\textit{unsupervised learning}), khususnya teknik klasterisasi pemartisian menggunakan algoritma K-Means dan K-Medoids. Kedua metode dievaluasi secara komparatif melalui dua studi kasus riil, yaitu dataset transaksi nasabah kartu kredit serta statistik lalu lintas penumpang penerbangan di Bandara Internasional San Francisco.

Algoritma K-Means terbukti unggul pada data yang memiliki distribusi kontinu dengan kepadatan kluster berbentuk bola (\textit{spherical cluster}). Keunggulan ini tercermin dari capaian Silhouette Score K-Means yang lebih tinggi dibandingkan K-Medoids pada kedua dataset uji, yakni 0,2795 berbanding 0,1945 pada dataset kartu kredit serta 0,4428 berbanding 0,3812 pada dataset lalu lintas penerbangan. Sebaliknya, algoritma K-Medoids memperlihatkan keunggulan representasi pusat kluster yang lebih realistis dan tangguh terhadap pengaruh nilai ekstrem atau pencilan (\textit{outliers}), mengingat penentuan medoid selalu merujuk pada sampel data aktual.

Kombinasi metode evaluasi internal berbasis kurva siku (\textit{Elbow Method}) menggunakan metrik Inertia dan validasi skor pemisahan siluet (\textit{Silhouette Score}) terbukti efektif dalam memandu penentuan jumlah kluster optimal secara terukur tanpa asumsi apriori. Analisis titik siku dengan bantuan pustaka \texttt{KneeLocator} memberikan rujukan kuantitatif yang konsisten dalam menetapkan konfigurasi partisi terbaik pada masing-masing eksperimen.

Penerapan klasterisasi pada kasus nyata membuka wawasan manajerial yang aplikatif. Segmentasi nasabah kartu kredit berhasil memetakan profil perilaku belanja dan kepatuhan pembayaran untuk perancangan promosi tertarget serta mitigasi risiko kredit. Sementara itu, pemetaan penumpang penerbangan bandara SFO menghasilkan tiga kelompok operasional yang jelas, sehingga manajemen dapat mengoptimalkan pemanfaatan gerbang keberangkatan, alokasi ruang terminal, dan fasilitas bagasi secara efisien.

%======================================================================
% DAFTAR PUSTAKA
%======================================================================
\cleardoublepage
\addcontentsline{toc}{chapter}{Daftar Pustaka}

\begin{thebibliography}{99}

\bibitem{afidaha2023}
Afidaha, M. W., \& Masrukan, M. H. (2023). Penerapan Metode Clustering dengan Algoritma K-Means untuk Pengelompokkan Data Migrasi Penduduk Tiap Kecamatan di Kabupaten Rembang. \textit{PRISMA, Prosiding Seminar Nasional Matematika} (pp. 729-738). Semarang: Jurusan Matematika, Universitas Negeri Semarang.

\bibitem{agrawal1993}
Agrawal, R., Imieliński, T., \& Swami, A. (1993). Mining association rules between sets of items in large databases. \textit{Proceedings of the 1993 ACM SIGMOD International Conference on Management of Data}, 207-216.

\bibitem{bishop2006}
Bishop, C. M. (2006). \textit{Pattern Recognition and Machine Learning}. Springer.

\bibitem{ester1996}
Ester, M., Kriegel, H.-P., Sander, J., \& Xu, X. (1996). A density-based algorithm for discovering clusters in large spatial databases with noise. \textit{Proceedings of the Second International Conference on Knowledge Discovery and Data Mining (KDD)}, 226-231.

\bibitem{han2011}
Han, J., Kamber, M., \& Pei, J. (2011). \textit{Data Mining: Concepts and Techniques} (3rd ed.). Morgan Kaufmann.

\bibitem{herviany2021}
Herviany, N., Delima, R., Nurhidayah, S., \& Kasini, A. (2021). Perbandingan Algoritma K-Means dan K-Medoids untuk Pengelompokkan Daerah Rawan Tanah Longsor di Provinsi Jawa Barat. \textit{MALCOM: Indonesian Journal of Machine Learning and Computer Science}, 34-40.

\bibitem{nurhalizah2024}
Nurhalizah, S., \& Ardianto, D. (2024). Penerapan Unsupervised Learning dalam Analisis Pola Data. \textit{Jurnal Sains Komputer dan Informatika}.

\bibitem{rousseeuw1987}
Rousseeuw, P. J. (1987). Silhouettes: A graphical aid to the interpretation and validation of cluster analysis. \textit{Journal of Computational and Applied Mathematics}, 20, 53-65.

\bibitem{sulistiyawati2021}
Sulistiyawati, A., \& Supriyanto, E. (2021). Pengelompokan Data Menggunakan Metode K-Means untuk Analisis Cluster. \textit{Jurnal Matematika dan Statistika}.

\bibitem{tan2019}
Tan, P.-N., Steinbach, M., \& Kumar, V. (2019). \textit{Introduction to Data Mining} (2nd ed.). Pearson.

\bibitem{wijoyo2024}
Wijoyo, S., et al. (2024). Unsupervised Learning and Its Application in Pattern Discovery. \textit{Jurnal Teknologi Informasi dan Sains Data}.

\bibitem{scikit2024}
scikit-learn. (2024). \textit{Clustering Performance Evaluation}. Dokumentasi resmi scikit-learn. Inertia \& Silhouette Score.

\end{thebibliography}

\end{document}
